\section{KXRCF troubled-cell indication}
\label{chap:kxrcf}

The KXRCF detector was introduced by Krivodonova, Xin, Remacle,
Chevaugeon, and Flaherty for discontinuous Galerkin approximations of
hyperbolic conservation laws~\cite{KrivodonovaEtAl2004}.  In contrast to a
limiter or to OFDG, a troubled-cell indicator does not modify the numerical
solution.  It supplies a binary decision about where a separate stabilization
mechanism should act.  This distinction is important in the present method:
KXRCF selects elements and the adapted OFDG operator defined in
\cref{chap:ofdg} damps their nonconstant polynomial content.

\subsection{Mathematical motivation}

\subsubsection{Interface superconvergence and jump detection}

Let $K\in\mathcal{T}_h$ be an element with outward unit normal
$\boldsymbol n_K$.  The DG approximation has independent traces on the two
sides of an interior face.  If $N$ denotes the neighboring element, define the
trace jump viewed from $K$ by
\begin{equation}
  [u_h]_{K,N}=u_h|_K-u_h|_N.
  \label{eq:kxrcf-jump}
\end{equation}
For a smooth exact solution, the two DG traces approximate the same physical
value and their jump decreases rapidly under mesh refinement.  The original
KXRCF construction exploits the particularly accurate DG trace on an outflow
boundary: at an inflow face of $K$, the neighboring trace is an outflow trace
of the upstream element.  A jump that is too large relative to the expected
smooth-mesh behavior therefore signals that at least one adjacent polynomial
does not represent a smooth solution.  This is the original motivation; the present preview does not measure superconvergence~\cite{KrivodonovaEtAl2004,QiuShu2005}.

Let $\boldsymbol v$ denote the transport velocity used solely for the
inflow/outflow decision.  The inflow part of the element boundary is
\begin{equation}
  \partial K^-
  =\{\boldsymbol x\in\partial K:
       \boldsymbol v(\boldsymbol x)\cdot\boldsymbol n_K(\boldsymbol x)<0\}.
  \label{eq:kxrcf-inflow}
\end{equation}
If there is no inflow measure, the indicator is zero. In the code, normal transport within $128\epsilon$ times the local speed scale is treated as tangential; $\epsilon$ is machine precision.  This
directional definition differs from sensors that inspect all faces
symmetrically: KXRCF deliberately measures information entering $K$.

\subsubsection{Normalized indicator}

Let $k$ be the reference-space order, $h_K$ the detector radius (distinct from
the OFDG volume-based length), and $S_K$ a scale of
the solution on $K$. Let $|\partial K^-|$ be the measure of the inflow boundary.
For a scalar DG solution, the KXRCF indicator used in this work is
\begin{equation}
 I_K=
 \frac{\left|\displaystyle\int_{\partial K^-}
        [u_h]_{K,N}\,\mathrm ds\right|}
      {h_K^{(k+1)/2}|\partial K^-|S_K}.
 \label{eq:kxrcf-indicator}
\end{equation}
The element is marked troubled when
\begin{equation}
  I_K>\tau,
  \qquad \tau=1 \text{ by default}.
  \label{eq:kxrcf-threshold}
\end{equation}
Here $\tau$ denotes the detector threshold, not the damping duration.
The face measure turns the numerator into an averaged signed jump, $S_K$
removes the local solution magnitude, and the factor $h_K^{(k+1)/2}$ separates
the expected smooth and discontinuous regimes.  If the face-averaged jump is
$\mathcal{O}(h_K^{k+1})$ in a smooth region and $S_K=\mathcal{O}(1)$, then
\begin{equation}
  I_K=\mathcal{O}\!\left(h_K^{(k+1)/2}\right)
  \longrightarrow 0.
  \label{eq:kxrcf-smooth-scaling}
\end{equation}
If an unresolved discontinuity produces an $\mathcal{O}(1)$ jump, the same
normalization instead gives
\begin{equation}
  I_K=\mathcal{O}\!\left(h_K^{-(k+1)/2}\right),
  \label{eq:kxrcf-shock-scaling}
\end{equation}
which grows under refinement.  These estimates explain the mesh-independent
threshold concept; they are an asymptotic motivation, not a new proof for the
modified multicomponent detector implemented here.  Comparative studies show
that detector behavior can still depend on polynomial degree, the chosen
indicator variable, and the threshold~\cite{QiuShu2005,FuShu2017,
VuikRyan2014,VuikRyan2016}.

\subsubsection{One-dimensional specialization}

In one dimension, $\partial K^-$ is one endpoint whenever the transport
velocity has a definite sign.  For $K=(x_{j-1/2},x_{j+1/2})$, positive
velocity, and $h_K=\Delta x_j/2$, \cref{eq:kxrcf-indicator} becomes
\begin{equation}
 I_j=
 \frac{\left|u_h(x_{j-1/2}^{+})-
                   u_h(x_{j-1/2}^{-})\right|}
      {(\Delta x_j/2)^{(k+1)/2}|\overline u_j|}.
 \label{eq:kxrcf-1d}
\end{equation}
The code uses the absolute cell average for the one-dimensional scale, matching
the common presentation of the detector.  In two and three dimensions it uses
the maximum absolute value at the element quadrature points.  The different
normalizations follow the original multidimensional use of a local maximum
norm rather than treating point boundaries as positive-measure faces.

\subsubsection{Signed integration and cancellation}

The absolute value in \cref{eq:kxrcf-indicator} is applied after the signed
jump has been integrated over all inflow faces.  Thus
\begin{equation}
 \left|\int_{\partial K^-}[u_h]_{K,N}\,\mathrm ds\right|
 \neq
 \int_{\partial K^-}|[u_h]_{K,N}|\,\mathrm ds
 \quad\text{in general}.
\end{equation}
The first expression is the KXRCF construction and permits cancellation along
or between inflow faces.  The second is an $L^1$ jump sensor of the type used by
the 2024 OEDG comparator.  Keeping these definitions distinct prevents the
implementation from silently turning KXRCF into a different detector.

\subsection{Hyperbolic systems and project-specific choices}

For component $c$ of a conservative state, the implementation forms
\begin{equation}
 I_{K,c}=
 \frac{\left|\displaystyle\int_{\partial K^-}
       (u_{K,c}-u_{N,c})\,\mathrm ds\right|}
      {h_K^{(k+1)/2}|\partial K^-|S_{K,c}},
 \qquad
 I_K=\max_c I_{K,c}.
 \label{eq:kxrcf-components}
\end{equation}
One active set is then shared by the complete conservative state,
\begin{equation}
  \mathcal A=\{K\in\mathcal T_h:I_K>\tau\}.
  \label{eq:kxrcf-active-set}
\end{equation}
This pooling prevents density, momentum, and energy from being filtered on
different cells.  It is a project-specific extension: published Euler uses of
KXRCF commonly monitor selected variables such as density, entropy, or total
energy rather than the maximum over every conservative component
\cite{KrivodonovaEtAl2004,FuShu2017}.

For scalar advection, the prescribed velocity determines inflow.  Burgers'
equation uses the secant velocity
\begin{equation}
  v_n=\tfrac12(u_1+u_2)n\quad\hbox{in one dimension},
\end{equation}
and Euler uses the arithmetic mean of the two fluid velocities,
\begin{equation}
  v_n=\tfrac12(\boldsymbol v_1+\boldsymbol v_2)
      \cdot\boldsymbol n.
  \label{eq:kxrcf-euler-inflow}
\end{equation}
The latter is an advective direction, not the Euler spectral radius.  Subsonic
Euler interfaces contain characteristic fields traveling in both directions,
so no single scalar velocity represents every characteristic family.  The damping methods therefore use the full one-sided radii
$|\boldsymbol v_i\cdot\boldsymbol n|+c_i$ for damping strength, evaluated from traces for adapted OFDG and element means for OEDG, whereas
\cref{eq:kxrcf-euler-inflow} is used only to decide which element receives the
KXRCF inflow contribution.  A characteristic-variable KXRCF detector remains
outside the present scope.

\subsection{Coupling to adapted OFDG}

After a completed Runge--Kutta step, KXRCF computes the active set
\cref{eq:kxrcf-active-set}.  A face is evaluated whenever at least one adjacent
element is active, ensuring that every active element receives all of its face
contributions.  The exact OFDG decay is then applied only to active elements.
Because the OFDG decay leaves the constant polynomial mode unchanged, this
gating does not alter cell-average conservation.  KXRCF therefore changes where
damping is applied, but neither its modal form nor its conservation argument.

\subsection{Implementation decisions}

\subsubsection{Geometry, quadrature, and normalization}

In one dimension $h_K$ is half the physical element length.  In higher
dimensions the implementation maps the reference-element center and vertices
and takes the largest sampled distance. Curved cells additionally include physical boundary-quadrature points; the resulting radius is a geometry-dependent sample bound, not an exact curved circumradius.  For affine symmetric cells this is their circumradius and a documented generalized
radius for other affine MFEM element families.  Volume quadrature supplies
$S_{K,c}$; face quadrature integrates the signed trace jump and the physical
inflow measure.  Each element and each side of a face owns its own evaluation
matrix, so supported meshes may mix finite-element signatures at one polynomial
order.  The verified cases combine triangles with quadrilaterals and combine
tetrahedra, hexahedra, and prisms.

A relative scale floor prevents division by a vanishing component scale.  It is
defined from the global maximum of the component scales and is used only when
the local scale is smaller.  A globally constant zero component has zero jump
and can never mark an element.  Physical boundary faces are omitted because no
neighboring DG trace exists; boundary conditions remain the responsibility of
the hyperbolic DG operator.

\subsubsection{Distributed ownership and caches}

Element evaluation matrices, quadrature weights, radii, local-face trace
matrices, and normals are immutable after construction.  During an indicator
application, ordinary interior faces update both adjacent elements.  An MPI
shared face updates only the locally owned element after MFEM exchanges the
neighbor coefficients; the neighboring rank performs the complementary
update.  All ranks participate in the component-scale reductions before any
root-only diagnostics are printed.

The standalone benchmark times complete \texttt{Compute} calls after warm-up.
No clocks, counters, or phase-specific timing branches are embedded in the
indicator itself; this keeps the core implementation focused on the numerical
method. The preview does not establish detector runtime savings.

\subsection{Limitations}

The implementation supports affine and static polynomial curvilinear, full-dimensional MFEM elements, with uniform solution order and a fixed mesh during the lifetime of the indicator. Mixed
geometry support is verified in two and three dimensions, including MPI faces
whose sides have different finite-element signatures. The supported geometry families and exclusions are listed in \cref{sec:support-boundary}. Its radius generalization,
component pooling, scale floor, boundary omission, and Euler inflow velocity
are engineering choices rather than claims from the original KXRCF analysis.
Moreover, selecting fewer cells does not itself guarantee that the resulting
OFDG--KXRCF method is nonoscillatory or positivity preserving.  These questions
are only partly addressed by the preview; detector selectivity remains the
explicit requirement KX-01 in \cref{chap:discussion}.
